% An IEEE conference paper (IEEEtran). For a journal, use [journal] instead of [conference].
\documentclass[conference]{IEEEtran}
\usepackage{amsmath,amssymb}
\usepackage{graphicx}
\usepackage{cite}

\begin{document}

\title{Title of the Paper}

\author{\IEEEauthorblockN{First Author}
\IEEEauthorblockA{Department, University\\
City, Country\\
first.author@example.org}
\and
\IEEEauthorblockN{Second Author}
\IEEEauthorblockA{Company\\
City, Country\\
second.author@example.org}}

\maketitle

\begin{abstract}
One paragraph: the problem, what you did, and what you found.
\end{abstract}

\begin{IEEEkeywords}
first keyword, second keyword, third keyword
\end{IEEEkeywords}

\section{Introduction}
Say why the problem matters and what this paper adds~\cite{knuth1984}.

\section{Method}
Describe the method. Number the equations you refer to:
\begin{equation}
  y = Wx + b. \label{eq:model}
\end{equation}
Equation~\eqref{eq:model} is the whole model.

\section{Results}
Paste or drop a figure into the text to add it here.

\section{Conclusion}
What the results mean, and what comes next.

\bibliographystyle{IEEEtran}
\bibliography{refs}

\end{document}
